\subsection{2.3 评测方法论}\label{evaluation-methodology}

我们使用一组专注于 NLL（负对数似然）任务的内部评测基准，用于实验与进度跟踪。它们来自不同的来源：供应商创建、与训练完全隔离的数据；互联网上不存在的内部来源；以及我们小心地将其从训练数据中移除的网络来源。我们的评测基准示例列于表 3，并在第 2.3.2 节讨论采用 NLL 评测基准的动机。我们将在第 2.7 节报告预训练后的评测结果。

我们总共拥有近 40 个此类评测基准，分为 5 个不同的类别：代码（Code）、STEM、数学（Math）\footnote{我们有意上调数学的权重，并将其与更宽泛的 STEM 类别区别对待。}、通用知识（General Knowledge）和多语言（Multilingual）。当我们需要将这些基准聚合为单一数值，以直接进行 EG 比较（第 2.2.2 节）或对模型排序时，我们使用以下公式：



Target = 0.5$\times$Coding+0.175$\times$STEM+0.175$\times$Math+0.1$\times$General
knowledge+0.05$\times$Multilingual (3)

这些显式权重使我们能够固化关于各评测基准相对优先级与重要性的决策。我们对代码和推理任务下游性能的重视，反映在预训练中我们为这些类别分配的更高权重上。

为计算式 3 中各类别得分，我们首先使用固定的内部参考模型作为基线对原始 NLL 结果进行归一化，然后对该类别中所有评测基准取均匀平均。所有被评测模型使用相同的 tokenizer，因此 NLL 值可直接比较。请注意，我们选择将数学提升为独立类别，而 STEM 则聚焦于科学、技术与工程。

表 3. 我们 NLL 评测套件中的数据示例类型。所有示例均经过仔细去重，确保与任何训练数据不重叠。

{\def\LTcaptype{none} % do not increment counter
\begin{longtable}[]{@{}lll@{}}
\toprule\noalign{}
\endhead
\bottomrule\noalign{}
\endlastfoot
类别 & 评测数据 & 来源 \\
\multirow{2}{*}{代码} & Microsoft 代码与 pull request & 内部私有项目 \\
& 人机协作编码会话 & 对先前模型的内部测试 \\
STEM & 研究生级别 STEM 问题的完整解答 & 由供应商定制 \\
数学 & 高等数学问题的完整解答 & 由供应商定制 \\
\multirow{4}{*}{通用知识} & 在线社区讨论 & 公共网络论坛 \\
& 人机交互 & 对先前模型的内部测试 \\
& Pyramidal trivia（金字塔式问答） & 与训练数据去重后的公开数据库 \\
& 高难度冷知识问答 & 由供应商定制 \\
多语言 & 多语言人机交互 & 对先前模型的内部测试 \\
\end{longtable}
}

\subsection{2.3.1 公共评测集净化}\label{public-evaluation-decontamination}

公开可用的评测基准在指导我们许多模型的开发方面发挥了重要作用。近年来，随着 AI 模型重要性的提升，评测数据集泄漏到训练数据来源中的情况也在增加。这在 GitHub 上尤为常见，许多 AI 研究人员和开发者会在 GitHub 上分享他们的工作。许多常见评测以其原始形式发布在 GitHub 上（Austin et al., 2021），几乎每个基准都会在那里被复制多次，通常附带解答、模型生成结果等内容。许多公共评测源自我们用于训练的部分在线来源，但经过了额外的整理流程（例如（Wei et al., 2024a; Hendrycks et al., 2020））。随着我们训练集的规模与覆盖面不断增长，这种重叠的影响变得显著。与先前工作类似（Godey et al., 2025），我们在实验中观察到，这种污染可能导致反直觉的结果，例如名义上旨在提升编码性能的数据同时也带来了长尾通用知识评测的提升。

我们采用了多种方法来检测并移除训练数据中泄漏的评测样本。所有来自 huggingface.co 及其镜像域名的数据均被移除，并对我们所有训练来源应用通用的 20-gram 模糊去重，相似度阈值为 80\%。我们发现这些方法能准确移除大量泄漏样本，但它们并不完美。



鉴于这一点，我们开发了仅供内部使用的私有基准，我们确信这些基准在互联网上的其他地方无法找到。

\subsection{2.3.2 准确率与 NLL 评测的比较}\label{comparison-of-accuracy-and-nll-evaluations}

我们将内部专有评测用于日常开发。这些评测由 MAI 团队成员策划，其中一部分由聘用领域专家的数据供应商定制，另一部分则收集自 Microsoft 全球组织内的领域专家。

对于预训练，我们选择聚焦于采用 NLL 评分方法的评测，而非选择题、生成式或其他形式。衡量模型能正确回答多少问题的评测在跨模型比较以及与人类基准对比时很方便，但它们存在下述若干操作层面的挑战。

评测的效率与成本。许多前沿基准需要思维链推理（Phan et al., 2025; Rein et al., 2023）、工具使用或其他形式的生成。这种自回归生成既昂贵又耗时。训练框架与配置往往并非为适合此类评测的快速自回归推理而优化。使用更高效的推理框架可以加速这些评测，但会引入额外的复杂度以及对架构变体的支持需求。采用开放式或自然语言答案的评测通常需要裁判模型（judge model），进一步增加成本。相比之下，NLL 评测与预训练具有相同的基本下一 token 预测目标，可以非常高效地运行。通过降低每次评测在时间与算力上的成本，我们可以运行更全面的评测，并确保每次实验之间的一致性。

对混杂因素的敏感性。评测可能对预训练模型行为中看似微小的变化高度敏感。例如，MMLU（Hendrycks et al., 2020）、MMLU-Pro（Wang et al., 2024b）及相关基准（Yue et al., 2024, 2025; Singh et al., 2025b）对评估预训练知识很有价值，但其选择题格式隐含地假设模型已具备解读和作答选择题的能力——这种能力在预训练评估中通常只有在规模出乎意料地大时才会出现（Du et al., 2024）。尽管这些任务有时采用基于 A、B、C、D 各选项 logits 的 NLL 评分方案，但许多问题需要的不仅仅是事实回忆，自由形式生成往往是一种更忠实的评估协议。

类似地，MATH（Hendrycks et al., 2021）要求严格遵守格式约定，例如将答案包裹在 \boxed{} 中；MBPP（Austin et al., 2021）则使模型暴露于共享 prompt 中的 Unix 风格换行符（\n）与问题陈述中 Windows 风格换行符（\r\n）之间的不一致。预训练数据中微小的格式差异——尤其是跨来源各不相同的情况——会显著影响模型对这些不一致的响应方式。我们在许多实验中观察到了此类交互，这使结果解读变得复杂，也增加了新老研究人员误读结果的可能性。

NLL 评测对此类问题更具鲁棒性。在每一步预测中，模型都以真实前缀为条件。这种教师强制（teacher-forcing）设置限制了微小误差不断累积的程度。我们的结论与先前工作一致，即基于损失或困惑度的指标能够在带有格式选择的任务上提供更高信噪比的模型开发反馈（Heineman et al., 2025; Du et al., 2024）。

构建的复杂性。构建一个高质量的问答（Q\&A）评测需要大量且精心协调的努力。通常需要多轮数据收集、难度校准、去重和质量控制，评测才能达到可接受的标准，且每个阶段都必须与领域专家密切协作完成。设计全新问题本身就是一项不平凡的挑战，尤其是在领域知识的前沿。



\pandocbounded{\includegraphics[keepaspectratio,alt={image}]{images/f6956ba66f001609e79548bfade2750671e1f400aa6ca269ecd56c9ca78338c5.jpg}}

图 4. 处理 HTML 预训练数据的流水线。我们从自研爬虫抓取和 Common Crawl 出发，先进行 HTML 提取，再依据各种启发式规则进行过滤。随后执行精确与模糊去重、跨来源去重以及 embedding 生成。这些数据随后由不同的流水线进一步处理；各流水线的具体细节见附录 B。

相比之下，基于 NLL 的评测的进入门槛要低得多。任何与主题相关的内容都可以作为初始语料，后续改进——例如施加额外结构、剪枝到更有代表性的样本，或纳入更高级的材料——会直接转化为评测质量的渐进提升。

\subsection{2.4 预训练数据}\label{pretraining-data}

MAI-Base-1 在公开可获取与获得许可的人类生成数据的混合语料上训练。我们汇集了涵盖网页、公开 GitHub 代码、图书、学术论文、新闻、多语言文本以及领域特定材料的高质量、多样化预训练数据。我们选择不在预训练中使用语言模型生成的合成数据，并努力避免和移除所收集数据来源中的 AI 生成内容（附录 B）。

我们所有的预训练数据均在内部处理，源自我们的几类基础来源：网络上的 HTML 与 PDF、图书与期刊，以及公开 GitHub 代码。我们不使用任何开源训练数据集，并将常见机器学习仓库排除在我们的网络数据之外，例如 huggingface.co 及类似网站。我们反复发现，对数据质量（而不仅仅是数量）进行仔细迭代至关重要——这不仅影响预训练结果，也影响推理基准上的下游表现。以下小节描述我们在处理各主要数据来源时采用的通用技术。每个来源数据流水线的具体细节与完整技术说明见附录 B。我们网络数据处理流水线的总览如图 4 所示。



表 4. 各训练数据来源的知识截止日期。

{\def\LTcaptype{none} % do not increment counter
\begin{longtable}[]{@{}|l|l|@{}}
\toprule\noalign{}
\endhead
\bottomrule\noalign{}
\endlastfoot
\hline
来源族系 & 知识截止日期 \\
\hline
网页 HTML & September 2025 \\
\hline
网页 PDF & December 2025 \\
\hline
公开 GitHub 代码 & June 2025 \\
\hline
图书与期刊 & March 2026 \\
\hline
\end{longtable}
}

\subsection{2.4.1 数据来源、许可与治理}\label{data-sources-licensing-and-governance}

我们数据中公开可获取的部分按照适用的使用条款和网络控制的行业标准进行收集。对于网络数据，我们使用尊重 Robots 排除协议（robots.txt）及相关 meta 标签和 HTML 控制的自研爬虫，使站点所有者能够管理其站点内容的访问与使用方式。我们还排除了违反 Microsoft 负责任 AI 政策或出现在美国贸易代表办公室（USTR）不良市场名单上的来源。

我们的训练语料还包括通过商业协议从第三方提供商处获取的数据集。第三方数据集需经过尽职调查流程，以评估数据保真度、所有权和适用使用权，并受保密条款与条件的约束。各来源的知识截止日期见表 4。

除用户明确选择加入或相关协议涵盖且尊重用户退出选择的情况外，我们未将私有客户数据或 Microsoft 产品与服务的数据用于预训练。训练之前，我们的整个语料都要经过 PII 风险与安全过滤；这些步骤的细节见附录 B。出于隐私、法律、安全与竞争方面的原因，我们不会披露完整的数据集清单或数据提供商。

\subsection{2.4.2 HTML 提取}\label{html-extraction}

我们很大一部分数据来自网络，其中可训练内容需要从原始 HTML 中提取。实现这一目标的方式有很多，所得质量差异往往相当细微。我们最感兴趣的部分内容——数学、代码等——在网络上的表示方式并不一致，许多现成方法最终会丢弃大量此类关键数据，同时保留周围的文本。表格和其他结构化内容也常常相当重要，但保留这些数据需要一些技巧。

我们不对所有来源使用单一的提取方法；相反，我们根据每个域名的结构和可靠性应用不同的方法：

• 针对特定来源的结构化解析器用于 HTML 或 XML 等标准化格式，基于 schema 的解析能够准确转换为文本表示。

• 使用 BeautifulSoup（Richardson, 2026）等工具实现的手工构造提取器，应用于结构一致但通用启发式规则处理不佳的域名，或内容价值与体量足以支撑额外工程投入的域名。

• LLM 与基于 agent 的处理用于需要定向提取、归一化或超出确定性规则能力的语义过滤的域名（Mahabadi et al., 2025）。在所有采用此方法的场景中，LLM 只能选择保留或删除原始来源文本中的内容，无法向输出中添加额外的合成内容。



• 当进一步处理有可能丢弃重要信息时，我们直接对原始内容进行训练。例如，Wikipedia 特有的标记语言 wikitext（MediaWiki, 2026）相比 HTML 缺乏成熟的工具链，而 infobox 等元素并不总能被现有解析器忠实保留。尽管 wikitext 的冗余度约为 3$\times$，我们发现直接训练完整的原始标记比依赖剥离简化后的转换效果更好。对于 Wikipedia 这类高价值、低体量的来源，这种取舍是值得的。

\subsection{2.4.3 去重}\label{deduplication}

我们这种规模与稀疏度的模型拥有非凡的知识获取与回忆能力。虽然这种能力是其作为通用系统发挥有效性的核心，但也增加了对训练数据的记忆与过拟合风险（Lee et al., 2022; Abbas et al., 2023）。当记忆发生时，模型无法泛化底层概念，导致在其他领域性能下降。此类过拟合的一个主要驱动因素是训练语料中的重复。

大规模网络数据包含大量冗余：许多文档互为近似重复，仅存在细微的格式差异，或在多个站点间被复制。若不进行仔细过滤，模型可能会接触到大量相同内容。随着模型容量的增加，这种冗余的影响变得更加显著。特别是，预测扩展行为对有效唯一 token 的数量敏感：当语料新颖性降低时，较大模型会表现出退化的扩展行为，因为它们更早地耗尽了新信息的供给。一旦数据集中的唯一内容饱和，额外的训练 token 几乎不再带来收益，从而限制了进一步的改进。在许多实验中，我们一致观察到，严格去重和移除冗余数据能同时改善预训练指标和下游推理基准。

鉴于限制冗余的重要性，我们的流水线包含多个去重阶段，每个阶段针对不同类别的冗余。我们的所有数据——无论来源如何——都会接受其中一种或多种去重处理。各来源（网络 HTML 与 PDF、图书与期刊、公开 GitHub 代码）的具体处理方式见附录 B 中的完整技术细节。这里我们介绍所采用的各种技术：

• 样板内容去除。网页 HTML 常包含重复的样板文本。利用文档内和跨文档的行出现统计，我们移除页眉、导航栏、页脚、侧边栏以及解析产物引入的冗余段落。

• 精确重复。相同内容可能因跨网站转载或因系统级故障造成的重复快照而多次出现。我们移除所有字节级与哈希级的精确重复。

• 模糊重复。我们应用基于 MinHash 局部敏感哈希（LSH）的模糊去重（Broder, 1997），相似度阈值为 0.8，遵循 Smith et al.~(2022) 的配置。

• 模板化网页。许多网站由共享模板生成页面，仅存在少量词汇变化（例如带原始算术表的“calculator”网页）。我们将每个页面骨架化为其最频繁出现的 token，并对这些模板执行模糊去重，以消除大批近似相同的页面。

• 语义重复。独立产出的文档由于共享上下文或规范的问题表述仍可能高度相似。这在代码数据集中尤为常见，众所周知的编程练习会在作业、考试、面试和竞赛中反复出现。尽管是独立编写的，解答往往趋于一致（例如二叉搜索树的标准遍历）。此类聚类在自然数据集中大量存在，移除它们能提升模型的鲁棒性和任务多样性（Sorscher et al., 2023; Abbas et al., 2023; Tirumala et al., 2023）。我们使用 Qwen3-Embedding-0.6B（Yang et al., 2025）模型识别语义相似的文档，并仅保留每个聚类中有限数量的代表样本。



我们既在每条处理流水线的早期阶段应用这些去重技术，也会在训练前作为最终步骤跨数据集再次应用。历史上，我们通过显式过滤器管理训练来源之间的重叠——例如，从所有网络爬取数据中排除 Wikipedia 域名。然而，随着我们来源的多样性与覆盖面不断增加，以及处理流水线变得更加复杂，跨数据集意外重叠的机会也在增长。

为解决这一问题，我们使用全局丢弃顺序（drop-order）在数据集上执行跨数据集去重。当检测到模糊或语义重复时，该实例仅保留在其出现的数据集中排名最高者中，并从所有丢弃顺序优先级更低的数据集中移除。因此，在解读实验结果和评估单个数据集的贡献或价值时，丢弃顺序的选择至关重要。对某个数据集的修改可能引入新的重叠，即使在未添加任何新内容的情况下，也会将数据从一个来源实际转移到另一个来源。

\subsection{2.4.4 数据过滤与分类}\label{filtering-and-categorization-of-data}

下游处理的一个核心目标是将庞大而异构的原始语料转化为既有利于预训练、又可受控地用于数据混合优化（第 2.5 节）的数据。在各来源中，我们首先移除不太可能对训练产生正向贡献的数据，包括垃圾内容、受限或政策敏感内容，以及其他来源特定的噪声。随后，我们将剩余数据归类为可解释的桶，例如质量分层、语言组、主题类别、教育价值、教育水平、来源类型和领域特定子语料。

由于不同的语料展现出不同的信号与失败模式，精确的过滤与分类策略因来源而异。因此，我们应用并组合了以下几类技术：

• 利用元数据信号，例如域名、文件名、仓库元数据、PDF 创建者与生产者字段，以及文档级元数据。

• 来源特定的启发式规则，例如网络文本质量过滤器、OCR 伪影过滤器、面向 STEM 内容的数学感知过滤器，以及针对生成代码的基于路径或内容的过滤器。

• 学习式分类器，包括 fastText 式分类器（Joulin et al., 2016）和基于 embedding 的模型，用于语言、主题、教育价值、教育水平、质量及其他语义属性。

• 提示式 LLM，选择性地用于更高价值或更模糊的决策，例如章节级提取、质量评判和细微的主题标注。

• 人工探索与标注，用于识别失败模式、验证过滤精度、审计高影响来源类别，并为分类器和 LLM 裁判构建训练数据。

过滤与分类共同将处理后的来源语料转化为结构化的数据桶集合。过滤通过移除低价值或有问题的内容提升数据质量，而分类则将剩余语料组织为可用于数据混合优化的可解释维度。这一结构支持在来源族系、质量分层、主题、教育水平、语言和来源特定格式之间进行受控消融。各主要来源的进一步实现细节见附录 B。

\subsection{2.4.5 数据消融实验}\label{data-ablations}

每个处理后的预训练数据集都要经过严格的消融。我们采用两种数据消融方法：单来源消融和扩展阶梯消融。单来源消融是一种隔离方法，用于量化特定数据集的边际效用，并将其贡献归因于基于 NLL 的评测。我们将目标数据在混合中的权重上调至 50\%，并从零开始训练。对于扩展阶梯消融，我们使用第 2.2.1 节中描述的扩展阶梯，在全量混合中消融该数据，以预测某项数据整理决策在所有其他数据集的背景下、于最终目标规模和 token 范围内将表现如何。阶梯消融中执行数据源降采样，以模拟目标模型的多轮次（multi-epoch）特性。



图 5. 数据混合前沿的示例说明。我们在 61 种 Web HTML、代码（Code）及其他数据的多样化混合配置上，以相同 TPP 在 3 种不同模型规模下从零训练了 183 个模型。实线表示 Web HTML 与代码的验证损失性能前沿。位于该 Pareto 前沿上的模型是这两个组分占主导的混合配置，而偏离前沿的模型则包含更大比例的“其他”（other）子集。

\pandocbounded{\includegraphics[keepaspectratio,alt={image}]{images/2851616d4fb6020782e7507739cd58de1cc835a730da35e6f7a69f5465cb6c1e.jpg}}
